\documentclass[11pt,a4paper]{article}
\usepackage{ftm}

\title{Conjectura de Hodge\\[4pt]
       \large Caderno de Lacunas Formais — Programa FTM}
\author{Federated Trivalent Manifold Research Programme}
\date{2026}

\begin{document}
\maketitle
\tableofcontents
\bigskip

%% ---------------------------------------------------------------
\section{Enunciado oficial}
%% ---------------------------------------------------------------

\begin{definition}[Conjectura de Hodge]
Seja $X$ uma variedade algébrica projetiva não singular sobre $\mathbb C$.
Uma classe $\alpha\in H^{2p}(X,\mathbb Q)$ é dita de \emph{Hodge}
se pertence a $H^{p,p}(X)$ na decomposição de Hodge.
A conjectura afirma que toda classe de Hodge é combinação
$\mathbb Q$-linear de classes de ciclos algébricos:
\[
  \mathrm{Hdg}^p(X)
  \;=\; H^{2p}(X,\mathbb Q)\cap H^{p,p}(X)
  \;\stackrel{?}{=}\; \mathrm{cl}\bigl(\mathrm{CH}^p(X)\otimes\mathbb Q\bigr).
\]
\end{definition}

\begin{remark}
A conjectura exige que $X$ seja projetiva e não singular sobre $\mathbb C$.
Variedades apenas topológicas, $PL$ ou $C^\infty$ estão \emph{fora}
do domínio; nenhum resultado obtido nessa categoria implica Hodge.
\end{remark}

\bigskip

%% ---------------------------------------------------------------
\section{Programas de pesquisa e encaixe FTM}
%% ---------------------------------------------------------------

\begin{programs}

\item \textbf{Caso abeliano.}
Cattani–Deligne–Kaplan (1995) provaram que o locus de Hodge de uma
família de variedades abelianas é uma reunião contável de subvariedades
algébricas. Para variedades abelianas e superfícies $K3$,
as classes de Hodge são geradas por ciclos de dimensão 1,
provando a conjectura nesse caso particular.

\begin{gap}
A prova não se estende a variedades gerais de dimensão $\ge3$:
o locus de Hodge pode ser denso em famílias com variação de
estrutura Hodge não trivial. \OPEN
\end{gap}

\item \textbf{Motivos e conjecturas padrão de Grothendieck.}
As quatro conjecturas padrão ($A$–$D$) implicariam Hodge via
a existência de uma categoria abeliana de motivos pura.
Em particular, a conjectura de Lefschetz (tipo $B$) daria
acesso a aplicações de ciclos racionais.

\begin{gap}
As conjecturas padrão de Grothendieck permanecem abertas;
em especial a conjectura de Lefschetz para variedades gerais. \OPEN
\end{gap}

\item \textbf{Teoria da deformação e Noether–Lefschetz.}
A variação de estrutura Hodge ao longo de uma família controla
quais classes permanecem do tipo $(p,p)$.
O teorema de Noether–Lefschetz garante que, para $d\ge4$,
uma superfície geral de grau $d$ em $\mathbb P^3$ tem
$\mathrm{Pic}(X)=\mathbb Z$.

A esperança é que, para ``quase toda'' deformação, as classes
de Hodge extras correspondam a ciclos; mas isso não foi provado.

\item \textbf{Hodge $p$-ádico e comparação.}
A teoria de Fontaine fornece isomorfismos de comparação entre
cohomologias Betti, étale e cristalina, levando a uma versão
$p$-ádica da decomposição de Hodge (Tsuji, Faltings).
Implicações para a conjectura de Hodge clássica via
comparação são conjecturais.

\end{programs}

%% ---------------------------------------------------------------
\section{Objeto FTM para variedades projetivas}
%% ---------------------------------------------------------------

\begin{definition}[Objeto FTM algébrico-complexo]
Para $X/\mathbb C$ projetiva não singular, o objeto FTM é
\[
  F_X=(X,\Gamma_X,\mathcal F_X,\lambda_X),
\]
onde $\Gamma_X$ é um complexo celular dual a uma cobertura de Stein,
$\mathcal F_X$ o feixe estrutural holomorfo e $\lambda_X$ os dados
de descida holomorfa. O objeto deve estar em $\mathcal R=C^\omega$
(analítico complexo).
\end{definition}

%% ---------------------------------------------------------------
\section{Obrigações formais FTM para Hodge}
%% ---------------------------------------------------------------

\begin{obligation}[H-DESC]\label{ob:HDESC}
A hipercobertura FTM calcula corretamente a cohomologia de $X$:
\[
  H^k(X,\mathbb Q)\cong H^k_{\mathrm{FTM}}(F_X,\mathbb Q).
\]
\end{obligation}

\begin{obligation}[H-TYPE]\label{ob:HTYPE}
A descida respeita a decomposição de Hodge:
\[
  H^k_{\mathrm{FTM}}(F_X,\mathbb C)
  =\bigoplus_{p+q=k}H^{p,q}_{\mathrm{FTM}}(F_X).
\]
\end{obligation}

\begin{obligation}[H-CYCLE]\label{ob:HCYCLE}
Classes de ciclos locais compatíveis colam em ciclo algébrico global:
\[
  \mathrm{cl}:\mathrm{CH}^p(X)\otimes\mathbb Q
  \longrightarrow \mathrm{Hdg}^p(X),
\]
com compatibilidade de restrição e colagem FTM.
\end{obligation}

\begin{obligation}[H-SURJ]\label{ob:HSURJ}
O mapa de ciclos $\mathrm{cl}$ é sobrejetor,
isto é, toda classe de Hodge é realizada por um ciclo algébrico.
Esta é a conjectura de Hodge propriamente dita.
\end{obligation}

\begin{remark}
Mesmo H-DESC e H-TYPE provados não implicam Hodge:
o obstáculo essencial é H-SURJ, que é independente
da construção do formalismo local.
\end{remark}

\begin{gap}[Sobrejetividade do mapa de ciclos]
Não existe prova de H-SURJ para variedades de dimensão $\ge2$
além dos casos cobertos pelas conjecturas padrão parciais
e pelo resultado de Cattani–Deligne–Kaplan.
H-DESC e H-TYPE são metas intermediárias; H-SURJ é a lacuna central. \OPEN
\end{gap}

\MATURITY{M0. H-DESC, H-TYPE, H-CYCLE e H-SURJ todas ABERTAS.}

\end{document}
